% [原文 p.549]
\section{代表全忠实函子的对象之构造}

\takeaway{本节只做构造：固定正合全忠实函子 $F$，借极丰富线丛给出的闭嵌入 $j:M\into\PP^{N}$，在 $\PP^{N}\times X$ 上造出对象 $\cE'$，再证它来自子簇 $M\times X$，即 $\cE'=\Rf J_{*}\cE$；工具是对角的消解、Koszul 型 $n$-消解与卷积的同构唯一性。}

本节从正合全忠实函子 $F$ 出发，构造 $\Db(M\times X)$ 中的某个对象 $\cE$；下一节将证明函子 $F$ 与 $\Phi_{\cE}$ 同构。$\cE$ 的构造分若干步完成。我们首先考虑一个闭嵌入 $j:M\into\PP^{N}$，并构造 $\Db(\PP^{N}\times X)$ 中的某个对象 $\cE'$；然后证明 $\cE'$ 其实来自子簇 $M\times X$，也就是说，存在对象 $\cE\in\Db(M\times X)$ 使 $\cE'=\Rf J_{*}\cE$，其中 $J=(j\times\id)$ 是 $M\times X$ 在 $\PP^{N}\times X$ 中的闭嵌入。

我们选取 $M$ 上一个极丰富线丛 $\cL$，使对一切 $k>0$ 与一切 $i\neq0$ 都有 \mbox{$H^{i}(\cL^{k})=0$}，并用 $j$ 记由 $\cL$ 确定的 $M$ 在 $\PP^{N}$ 中的闭嵌入。

积 $\PP^{N}\times\PP^{N}$ 有所谓的对角消解（见 [2]）。它是形如
\begin{equation}
\begin{aligned}
0\longrightarrow{}&\cO(-N)\boxtimes\Omega^{N}(N)\xrightarrow{d_{-N}}\cO(-N+1)\boxtimes\Omega^{N-1}(N-1)\\
&\longrightarrow\cdots\longrightarrow\cO(-1)\boxtimes\Omega^{1}(1)\xrightarrow{d_{-1}}\cO\boxtimes\cO.
\end{aligned}
\tag{31}
\end{equation}
的层的复形。该复形是结构层 $\cO_{\Delta}$ 的消解，其中 $\Delta$ 是积 $\PP^{N}\times\PP^{N}$ 的对角。

\yizhu{式 (31) 就是 Beilinson 为 $\PP^{N}\times\PP^{N}$ 的对角所造的 Koszul 型消解：第 $-i$ 项取 $\cO(-i)\boxtimes\Omega^{i}(i)$，微分由外积与缩并的配对给出，整条复形只用到 $\PP^{N}$ 上的欧拉序列。把 $\Omega^{i}(i)$ 逐项换成 $\bigwedge^{i}A_{1}\otimes\cO$，就得到后文 (35) 中射影空间上的消解；两者是同一个 Koszul 复形的几何版与代数版。【来源：原书 p.549 式 (31) 及 p.553 式 (35)。】}

记 $F'$ 为从 $\Db(\PP^{N})$ 到 $\Db(X)$ 的函子，即复合 $F\circ\Lf j^{*}$；并考虑投影构成的图
\diag{\PP^{N}\times X \arrow[r,"\pi'"] \arrow[d,"q"'] \& X\\
\PP^{N} \&}

% [原文 p.550]
我们记
\begin{equation*}
d'_{-i}\in\Hom_{\PP^{N}\times X}\bigl(\cO(-i)\boxtimes F'(\Omega^{i}(i)),\ \cO(-i+1)\boxtimes F'(\Omega^{i-1}(i-1))\bigr)
\end{equation*}
为态射 $d_{-i}$ 在如下复合映射下的像：
\begin{align*}
&\Hom\bigl(\cO(-i)\boxtimes\Omega^{i}(i),\ \cO(-i+1)\boxtimes\Omega^{i-1}(i-1)\bigr)\\
&\qquad\iso\Hom\bigl(\cO\boxtimes\Omega^{i}(i),\ \cO(1)\boxtimes\Omega^{i-1}(i-1)\bigr)\\
&\qquad\iso\Hom\bigl(\Omega^{i}(i),\ H^{0}(\cO(1))\otimes\Omega^{i-1}(i-1)\bigr)\\
&\qquad\longrightarrow\Hom\bigl(F'(\Omega^{i}(i)),\ H^{0}(\cO(1))\otimes F'(\Omega^{i-1}(i-1))\bigr)\\
&\qquad\iso\Hom\bigl(\cO\boxtimes F'(\Omega^{i}(i)),\ \cO(1)\boxtimes F'(\Omega^{i-1}(i-1))\bigr)\\
&\qquad\iso\Hom\bigl(\cO(-i)\boxtimes F'(\Omega^{i}(i)),\ \cO(-i+1)\boxtimes F'(\Omega^{i-1}(i-1))\bigr).
\end{align*}

容易看出，复合 $d'_{-i+1}\circ d'_{-i}$ 为零。因此，我们可以考虑导出范畴 $\Db(\PP^{N}\times X)$ 中对象的如下有界复形：
\begin{equation}
C^{\bullet}:=\left\{\begin{aligned}
&\cO(-N)\boxtimes F'(\Omega^{N}(N))\xrightarrow{d'_{-N}}\cdots\\
&\qquad\longrightarrow\cO(-1)\boxtimes F'(\Omega^{1}(1))\xrightarrow{d'_{-1}}\cO\boxtimes F'(\cO)
\end{aligned}\right\}.
\tag{32}
\end{equation}

对 $l<0$ 有
\begin{align*}
&\Hom^{l}\bigl(\cO(-i)\boxtimes F'(\Omega^{i}(i)),\ \cO(-k)\boxtimes F'(\Omega^{k}(k))\bigr)\\
&\qquad\iso\Hom^{l}\bigl(\cO\boxtimes F'(\Omega^{i}(i)),\ H^{0}(\cO(i-k))\otimes F'(\Omega^{k}(k))\bigr)\\
&\qquad\iso\Hom^{l}\bigl(j^{*}(\Omega^{i}(i)),\ H^{0}(\cO(i-k))\otimes j^{*}(\Omega^{k}(k))\bigr)=0.
\end{align*}

于是由引理 3.1.2，$C^{\bullet}$ 有卷积，且所有卷积彼此同构。记 $\cE'$ 为 $C^{\bullet}$ 的一个卷积，$\gamma_{0}$ 为态射 $\cO\boxtimes F'(\cO)\xrightarrow{\gamma_{0}}\cE'$。（事实上，下文会看到 $C^{\bullet}$ 的所有卷积都典范同构。）现在令 $\Phi_{\cE'}$ 为 $\Db(\PP^{N})$ 到 $\Db(X)$ 的、由 (7) 定义的函子。

\yizhu{这里所谓复形 $C^{\bullet}$ 的卷积，是把各项按次数平移后拼成一个对象：取 $\bigoplus_{i}C^{i}[-i]$，微分取 $d_{C}$ 与平移带来的符号 $(-1)^{i}$ 之和；其存在即引理 3.1.2 的结论，该引理同时保证，当各项之间的负次数 $\Hom$ 全为零时这样的对象在同构意义下唯一。后文反复使用这一点：只要认出两个对象是同一个复形的卷积，就得到它们同构。【来源：原书 p.550 对引理 3.1.2 的援引，及 p.551、p.553 的同类论证。】}

\begin{lemma}{3.3.1}
对一切 $k\in\Z$ 都有典范同构
\begin{equation*}
f_{k}:F'(\cO(k))\iso\Phi_{\cE'}(\cO(k)),
\end{equation*}
并且这些同构是函子性的；也就是说，对任一 $\alpha:\cO(k)\to\cO(l)$，图
\diag{F'(\cO(k)) \arrow[r,"F'(\alpha)"] \arrow[d,"f_{k}"'] \& F'(\cO(l)) \arrow[d,"f_{l}"']\\
\Phi_{\cE'}(\cO(k)) \arrow[r,"\Phi_{\cE'}(\alpha)"] \& \Phi_{\cE'}(\cO(l))}
交换。
\end{lemma}

\begin{proof}
先设 $k\geqslant0$，考虑对角 $\Delta\subset\PP^{N}\times\PP^{N}$ 的消解 (31)。把它张量 $\cO(k)\boxtimes\cO$，再取它在第二个因子上的
% [原文 p.551]
投影下的直接像。结果得到射影空间 $\PP^{N}$ 上的层 $\cO(k)$ 的如下消解：
\begin{equation*}
\begin{aligned}
&\Bigl\{H^{0}(\cO(k-N))\otimes\Omega^{N}(N)\longrightarrow\cdots\\
&\qquad\longrightarrow H^{0}(\cO(k-1))\otimes\Omega^{1}(1)\longrightarrow H^{0}(\cO(k))\otimes\cO\Bigr\}\xrightarrow{\ \delta_{k}\ }\cO(k).
\end{aligned}
\end{equation*}
由于按假设 $F'$ 正合，可见 $F'(\cO(k))$ 是复形
\begin{equation*}
\begin{aligned}
&H^{0}(\cO(k-N))\otimes F'(\Omega^{N}(N))\longrightarrow\cdots\\
&\qquad\longrightarrow H^{0}(\cO(k-1))\otimes F'(\Omega^{1}(1))\longrightarrow H^{0}(\cO(k))\otimes F'(\cO)
\end{aligned}
\end{equation*}
的卷积，该复形的各项都是范畴 $\Db(X)$ 的对象。我们把这个复形记为 $D^{\bullet}_{k}$。

现在回忆，按构造，$\cE'$ 是复形 $C^{\bullet}$ (32) 的卷积。考虑 $\PP^{N}\times X$ 上的复形 $C^{\bullet}_{k}:=q^{*}\cO(k)\otimes C^{\bullet}$。那么 $q^{*}\cO(k)\otimes\cE'$ 是 $C^{\bullet}_{k}$ 的卷积；并且有由 $\gamma_{0}$ 典范得到的态射 $\gamma_{k}:\cO(k)\boxtimes F'(\cO)\to q^{*}\cO(k)\otimes\cE'$。复形 $\pi'_{*}(C^{\bullet}_{k})$，即 $(C^{\bullet}_{k})$ 在第二个因子投影下的直接像，典范同构于 $D^{\bullet}_{k}$。这样我们看到，对象 $F'(\cO(k))$ 与 $\Phi_{\cE'}(\cO(k)):=\Rf\pi'_{*}(q^{*}\cO(k)\otimes\cE')$ 都是同一个复形 $D^{\bullet}_{k}$ 的卷积。
% [Tutor Note: 页图此句印作 $\pi'_{*}(C^{\bullet}_{k})$，没有粗体 $\mathbf{R}$；而同一句后半与 p.554 以下一律印 $\mathbf{R}\pi'_{*}$。此处照录页图，读者按导出的直接像理解即可。]

按假设，函子 $F$ 全忠实。因此对 $\PP^{N}$ 上局部自由层 $\cG$ 与 $\cH$ 有等式
\begin{equation*}
\Hom^{i}(F'(\cG),F'(\cH))=\Hom^{i}(j^{*}(\cG),j^{*}(\cH))=0\qquad(i<0).
\end{equation*}
特别地，这说明复形 $D^{\bullet}_{k}$ 满足引理 3.1.2 的条件 (25) 与 (26)。于是由该引理，存在唯一确定的同构 $f_{k}:F'(\cO(k))\iso\Phi_{\cE'}(\cO(k))$，使下图交换：
\diag{H^{0}(\cO(k))\otimes F'(\cO) \arrow[r,"F'(\delta_{k})"] \arrow[d,"\id"'] \& F'(\cO(k)) \arrow[d,"f_{k}"]\\
H^{0}(\cO(k))\otimes F'(\cO) \arrow[r,"\Rf\pi'_{*}(\gamma_{k})"] \& \Phi_{\cE'}(\cO(k))}

下面证明这些同构是函子性的。对任一 $\alpha:\cO(k)\to\cO(l)$，有如下形式的交换方块
\diag{H^{0}(\cO(k))\otimes F'(\cO) \arrow[r,"F'(\delta_{k})"] \arrow[d,"H^{0}(\alpha)\otimes\id"'] \& F'(\cO(k)) \arrow[d,"F'(\alpha)"]\\
H^{0}(\cO(l))\otimes F'(\cO) \arrow[r,"F'(\delta_{l})"] \& F'(\cO(l))}
以及
\diag{H^{0}(\cO(k))\otimes F'(\cO) \arrow[r,"\Rf\pi'_{*}(\gamma_{k})"] \arrow[d,"H^{0}(\alpha)\otimes\id"'] \& \Phi_{\cE'}(\cO(k)) \arrow[d,"\Phi_{\cE'}(\alpha)"]\\
H^{0}(\cO(l))\otimes F'(\cO) \arrow[r,"\Rf\pi'_{*}(\gamma_{l})"] \& \Phi_{\cE'}(\cO(l))}
% [原文 p.552]
这三个交换方块蕴含如下等式：
\begin{align*}
f_{l}\circ F'(\alpha)\circ F'(\delta_{k})&=f_{l}\circ F'(\delta_{l})\circ(H^{0}(\alpha)\otimes\id)=\Rf\pi'_{*}(\gamma_{l})\circ(H^{0}(\alpha)\otimes\id),\\
\Phi_{\cE'}(\alpha)\circ f_{k}\circ F'(\delta_{k})&=\Phi_{\cE'}(\alpha)\circ\Rf\pi'_{*}(\gamma_{k})=\Rf\pi'_{*}(\gamma_{l})\circ(H^{0}(\alpha)\otimes\id).
\end{align*}
复形 $D^{\bullet}_{k}$ 与 $D^{\bullet}_{l}$ 满足引理 3.1.3 的条件，因此存在唯一态射 $h:F'(\cO(k))\to\Phi_{\cE'}(\cO(l))$，使
\begin{equation*}
h\circ F'(\delta_{k})=\Rf\pi'_{*}(\gamma_{l})\circ(H^{0}(\alpha)\otimes\id).
\end{equation*}
于是态射 $h$ 同时与 $f_{l}\circ F'(\alpha)$ 及 $\Phi_{\cE'}(\alpha)\circ f_{k}$ 重合，这就说明这两者相等。

现在考虑 $k<0$ 的情形。取 $\PP^{N}$ 上的层 $\cO(k)$ 的右消解
\begin{equation*}
\cO(k)\iso\Bigl\{V^{k}_{0}\otimes\cO\longrightarrow\cdots\longrightarrow V^{k}_{N}\otimes\cO(N)\Bigr\}
\end{equation*}
再次应用引理 3.1.3，可见复形映射
\begin{center}\begin{tikzcd}[row sep=2.4em,column sep=1.9em,ampersand replacement=\&]
V^{k}_{0}\otimes F'(\cO) \arrow[r] \arrow[d,"\id\otimes f_{0}"',"{\wr}"{right}] \& \cdots \& V^{k}_{N}\otimes F'(\cO(N)) \arrow[d,"\id\otimes f_{N}"',"{\wr}"{right}]\\
V^{k}_{0}\otimes\Phi_{\cE'}(\cO) \arrow[r] \& \cdots \& V^{k}_{N}\otimes\Phi_{\cE'}(\cO(N))
\end{tikzcd}\end{center}
给出唯一确定的态射 $f_{k}:F'(\cO(k))\to\Phi_{\cE'}(\cO(k))$。直接验证（我们略去）表明这些态射是函子性的。
\end{proof}

\begin{remark}{3.3.2}
我们注意到，由函子 $F$ 构造出的对象 $\cE'\in\Db(\PP^{N}\times X)$ 是唯一确定的。
\end{remark}

现在我们来证明：范畴 $\Db(M\times X)$ 中存在对象 $\cE$，使 $\Rf J_{*}\cE\cong\cE'$，其中如上所述，$J$ 是 $M\times X$ 在 $\PP^{N}\times X$ 中的嵌入。
% [Tutor Note: 页图此处印作 “the existence of an object in the category $E\in\mathbf{D}^{b}(M\times X)$”，把对象名 $E$ 与范畴名错位；页图与文字层一致，非抽字错误。译文按原意作“范畴 $\Db(M\times X)$ 中存在对象 $\cE$”。]

设 $\cL$ 是 $M$ 上的极丰富线丛，$j:M\into\PP^{N}$ 是它给出的到射影空间的嵌入。用 $A$ 记分次代数 $\bigoplus_{i=0}^{\infty}H^{0}(M,\cL^{i})$。

令 $B_{0}=k$，$B_{1}=A_{1}$。对 $m\geqslant2$，按如下规则定义 $B_{m}$：
\begin{equation}
B_{m}=\Ker\bigl(B_{m-1}\otimes A_{1}\xrightarrow{u_{m-1}}B_{m-2}\otimes A_{2}\bigr),
\tag{33}
\end{equation}
其中 $u_{m-1}$ 是由归纳法定义的自然映射。

\begin{definition}{3.3.3}
若右 $A$-模的序列
\begin{equation*}
B_{n}\otimes_{k}A\longrightarrow B_{n-1}\otimes_{k}A\longrightarrow\cdots\longrightarrow B_{1}\otimes_{k}A\longrightarrow A\longrightarrow k\longrightarrow0
\end{equation*}
正合，则称代数 $A$ 是 $n$-Koszul 代数。若一个代数是任一 $n$ 的 $n$-Koszul 代数，则称之为 Koszul 代数。
\end{definition}

\yizhu{$n$-Koszul 性是 Koszul 性的有限截断：只要求前 $n$ 步的模序列正合，不问整条消解。由 (33) 递归给出的 $B_{m}$ 正是 Koszul 复形的各项；p.553 的自然嵌入 $\bigwedge^{i}A_{1}\subset B_{i}$ 说明 $B_{i}$ 至少含外代数的线性部分，这正是 Koszul 性在交换情形下的含义。【来源：原书 p.552 定义 3.3.3、p.553 的嵌入，以及文献 [1]（Backelin）关于 Veronese 代数的结果。】}

设 $A$ 是 $n$-Koszul 代数。令 $R_{0}=\cO_{M}$；对 $m\geqslant1$，用 $R_{m}$ 记典范态射
\begin{equation*}
B_{m}\otimes\cO_{M}\longrightarrow B_{m-1}\otimes\cL
\end{equation*}
的核，该态射由自然嵌入 $B_{m}\into B_{m-1}\otimes A_{1}$ 定义。
% [原文 p.553]
利用 (33)，我们得到典范态射 $R_{m}\to A_{1}\otimes R_{m-1}$（事实上可以验证，有同构 $\Hom(R_{m},R_{m-1})\cong A_{1}^{*}$）。

此外，若 $A$ 是 $n$-Koszul 代数，则对 $m\leqslant n$，层的如下复形正合：
\begin{equation}
0\longrightarrow R_{m}\longrightarrow B_{m}\otimes\cO_{M}\longrightarrow B_{m-1}\otimes\cL\longrightarrow\cdots\longrightarrow B_{1}\otimes\cL^{m-1}\longrightarrow\cL^{m}\longrightarrow0.
\tag{34}
\end{equation}

射影空间 $\PP^{N}$ 上有形如
\begin{equation}
0\longrightarrow\Omega^{m}(m)\longrightarrow\bigwedge^{m}A_{1}\otimes\cO\longrightarrow\bigwedge^{m-1}A_{1}\otimes\cO(1)\longrightarrow\cdots\longrightarrow\cO(m)\longrightarrow0.
\tag{35}
\end{equation}
的正合复形。存在典范映射 $f_{m}:j^{*}\Omega^{m}(m)\to R_{m}$。确实，由于 $A$ 交换，有自然嵌入 $\bigwedge^{i}A_{1}\subset B_{i}$。因此存在从限制到 $M$ 的复形 (35) 到复形 (34) 的态射，从而有典范映射 $f_{m}:j^{*}\Omega^{m}(m)\to R_{m}$。

已知：对任一 $n$ 都存在 $l$，使 Veronese 代数 $A^{l}=\bigoplus_{i=0}^{\infty}H^{0}(M,\cL^{il})$ 是 $n$-Koszul 的；此外，[1] 中证明了当 $l\gg0$ 时代数 $A^{l}$ 事实上是 Koszul 代数。

然而在下文中，除 Veronese 代数的 $n$-Koszul 性质之外，我们还需要一些附加性质。也就是说，利用 [18] 的方法，把层 $\cL$ 换成它的足够高次幂 $\cL^{j}$，可以证明如下断言。

\begin{proposition}{3.3.4}
对任一整数 $n$，存在极丰富线丛 $\cL$，使得
\begin{enumerate}[label=\arabic*),leftmargin=2.6em,itemsep=2pt,topsep=2pt]
\item 代数 $A$ 是 $n$-Koszul 代数，也就是说，序列
\begin{equation*}
B_{n}\otimes_{k}A\longrightarrow B_{n-1}\otimes_{k}A\longrightarrow\cdots\longrightarrow B_{1}\otimes_{k}A\longrightarrow A\longrightarrow k\longrightarrow0
\end{equation*}
正合；
\item 由
\begin{align*}
&A_{k-n}\otimes R_{n}\longrightarrow A_{k-n+1}\otimes R_{n-1}\longrightarrow\\
&\qquad\cdots\longrightarrow A_{k-1}\otimes R_{1}\longrightarrow A_{k}\otimes R_{0}\longrightarrow\cL^{k}\longrightarrow0
\end{align*}
给出的 $M$ 上的层复形对任一 $k\geqslant0$ 正合（若 $k-i<0$，则按定义 $A_{k-i}=0$）；
\item 形如
\begin{equation*}
\cL^{-n}\boxtimes R_{n}\longrightarrow\cdots\longrightarrow\cL^{-1}\boxtimes R_{1}\longrightarrow\cO_{M}\boxtimes R_{0}\longrightarrow\cO_{\Delta}
\end{equation*}
的 $M\times M$ 上的层复形正合，也就是说，它给出 $M\times M$ 上对角的一个 $n$-消解。
\end{enumerate}
\end{proposition}

本命题的证明见 \S\,3.5。

用 $T_{k}$ 记典范态射 $A_{k-n}\otimes R_{n}\to A_{k-n+1}\otimes R_{n-1}$ 的核。鉴于命题 3.3.4 的性质 2) 以及当 $n\gg0$ 时 $\Ext^{n+1}(\cL^{k},T_{k})=0$ 这一事实，可见复形
\begin{equation*}
A_{k-n}\otimes R_{n}\longrightarrow A_{k-n+1}\otimes R_{n-1}\longrightarrow\cdots\longrightarrow A_{k}\otimes R_{0}
\end{equation*}
的每个卷积都典范同构于 $T_{k}[n]\oplus\cL^{k}$。

% [原文 p.554]
典范态射 $R_{k}\to A_{1}\otimes R_{k-1}$ 诱导出态射
\begin{equation*}
\cL^{-k}\boxtimes F(R_{k})\longrightarrow\cL^{-k+1}\boxtimes F(R_{k-1}).
\end{equation*}
这来自如下同构的存在：
\begin{align*}
\Hom\bigl(\cL^{-k}\boxtimes F(R_{k}),\ \cL^{-k+1}\boxtimes F(R_{k-1})\bigr)&\cong\Hom\bigl(F(R_{k}),\ H^{0}(\cL)\otimes F(R_{k-1})\bigr)\\
&\cong\Hom(R_{k},\ A_{1}\otimes R_{k-1}).
\end{align*}
此外，我们有范畴 $\Db(M\times X)$ 中对象的如下复形：
\begin{equation}
\cL^{-n}\boxtimes F(R_{n})\longrightarrow\cdots\longrightarrow\cL^{-1}\boxtimes F(R_{1})\longrightarrow\cO_{M}\boxtimes F(R_{0}).
\tag{36}
\end{equation}

由引理 3.1.2，复形 (36) 有卷积，且它的所有卷积彼此同构。我们把这个卷积记为 $G\in\Db(M\times X)$。

对任一 $k\geqslant0$，对象 $\Rf\pi_{*}(G\otimes p^{*}(\cL^{k}))$ 是复形
\begin{equation}
A_{k-n}\otimes F(R_{n})\longrightarrow A_{k-n+1}\otimes F(R_{n-1})\longrightarrow\cdots\longrightarrow A_{k}\otimes F(R_{0}).
\tag{37}
\end{equation}
的卷积。另一方面，对象 $F(T_{k}[n]\oplus\cL^{k})$ 也是该复形的卷积，显然满足引理 3.1.2 的条件。因此有同构 $\Rf\pi_{*}(G\otimes p^{*}(\cL^{k}))\cong F(T_{k}[n]\oplus\cL^{k})$。

由引理 3.2.3 与注 3.2.4，对一切 $k>0$，非平凡上同调层 $H^{i}(\Rf\pi_{*}(G\otimes p^{*}(\cL^{k})))=H^{i}(F(T_{k})[n])\oplus H^{i}(F(\cL^{k}))$ 集中在并集 $[-n-a,-n]\cup[-a,0]$ 中（其中 $a$ 是注 3.2.4 中定义的数）。由于 $\cL$ 丰富，上同调层 $H^{i}(G)$ 也集中在 $[-n-a,-n]\cup[-a,0]$ 中。我们可以假定 $n>\dim M+\dim X+a$。由于 $M\times X$ 上凝聚层范畴的同调维数等于 $\dim M+\dim X$，在这种情形下我们看到 $G\cong C\oplus\cE$，其中 $\cE$、$C$ 是 $\Db(M\times X)$ 的对象，且当 $i\notin[-a,0]$ 时 $H^{i}(\cE)=0$、当 $i\notin[-n-a,-n]$ 时 $H^{i}(C)=0$。于是特别地 $\Rf\pi_{*}(\cE\otimes p^{*}(\cL^{k}))\cong F(\cL^{k})$。注意，由于对象 $G$ 作为复形 (36) 的卷积唯一确定，对象 $\cE$ 也在同构意义下唯一确定。

下面证明存在同构 $\Rf J_{*}\cE\cong\cE'$。为此，考虑 $\Db(\PP^{N}\times X)$ 上的复形映射
\begin{center}\begin{tikzcd}[row sep=2.7em,column sep=2.2em,ampersand replacement=\&]
\cO(-n)\boxtimes F'(\Omega^{n}(n)) \arrow[r] \arrow[d,"\mathrm{can}\boxtimes F(f_{n})"] \& \cdots \& \cO\boxtimes F'(\cO) \arrow[d,"\mathrm{can}\boxtimes F(f_{0})"]\\
\Rf j_{*}\cL^{-n}\boxtimes F(R_{n}) \arrow[r] \& \cdots \& \Rf j_{*}\cO_{M}\boxtimes F(R_{0})
\end{tikzcd}\end{center}
% [Tutor Note: 页图此图下行的左端印作 $\mathbf{R}j_{*}\cL^{-n}$ 与 $\mathbf{R}j_{*}\cO_{M}$（小写 $j$），而 p.549、p.552、p.555 处同一个函子一律印 $\mathbf{R}J_{*}$，其中 $J=(j\times\id)$。据上下文此处应为 $\mathbf{R}J_{*}$；此处照录页图，不改字。]
应用引理 3.1.3，我们得到卷积之间的态射 $\varphi:K\to\Rf J_{*}G$ 的存在性。

若 $N>n$，则对象 $K$ 与 $\cE'$ 不同构，但有特异三角
\begin{equation*}
S\longrightarrow K\longrightarrow\cE'\longrightarrow S[1].
\end{equation*}

% [原文 p.555]
同上可以证明，上同调层 $H^{i}(S)$ 仅在 $i\in[-n-a,-n]$ 时非零。这说明 $\Hom(S,\Rf J_{*}\cE)=0$ 且 $\Hom(S[1],\Rf J_{*}\cE)=0$，因为上同调 $\Rf J_{*}\cE$ 集中在闭区间 $[-a,0]$ 内。由此可知存在唯一态射 $\psi:\cE'\to\Rf J_{*}\cE$，使下图交换：
\diag{K \arrow[r,"\varphi"] \arrow[d] \& \Rf J_{*}G \arrow[d]\\
\cE' \arrow[r,"\psi"] \& \Rf J_{*}\cE}

如我们所知，
\begin{equation*}
\Rf\pi'_{*}(\cE'\otimes q^{*}(\cO(k)))\cong F(\cL^{k})\cong\Rf\pi_{*}(\cE\otimes p^{*}(\cL^{k})).
\end{equation*}
用 $\psi_{k}$ 记由 $\psi$ 诱导的态射 $\Rf\pi'_{*}(\cE'\otimes q^{*}(\cO(k)))\to\Rf\pi_{*}(\cE\otimes p^{*}(\cL^{k}))$。这些 $\psi_{k}$ 装进交换图
\begin{center}\begin{tikzcd}[row sep=2.2em,column sep=1.3em,ampersand replacement=\&]
S^{k}A_{1}\otimes F(\cO) \arrow[r,"\mathrm{can}"] \arrow[d,"\mathrm{can}"'] \& F(\cL^{k}) \arrow[r,"\sim"] \& \Rf\pi'_{*}(\cE'\otimes q^{*}(\cO(k))) \arrow[d,"\psi_{k}"]\\
A_{k}\otimes F(\cO) \arrow[r,"\mathrm{can}"] \& F(\cL^{k}) \arrow[r,"\sim"] \& \Rf\pi_{*}(\cE\otimes p^{*}(\cL^{k}))
\end{tikzcd}\end{center}
% [Tutor Note: 页图此图左上印 $S^{k}A_{1}\otimes F(\cO)$，左下印 $A_{k}\otimes F(\cO)$，二者并置。这并不是笔误：当分次代数 $A$ 由 $A_{1}$ 生成时有 $S^{k}A_{1}=A_{k}$，此处照录页图。]

这意味着对任一 $k\geqslant0$，态射 $\psi_{k}$ 都是同构。因此 $\psi$ 也是同构。这样我们就证明了如下断言。

\begin{proposition}{3.3.5}
存在对象 $\cE\in\Db(M\times X)$，使 $\Rf J_{*}\cE\cong\cE'$，其中 $\cE'$ 是 \S\,3.3 中构造的 $\Db(\PP^{N}\times X)$ 的对象；并且这个 $\cE$ 在同构意义下唯一。
\end{proposition}
% [Tutor Note: 页图此条末尾 “unique up to isomorphism” 之后无句点，与全书陈述末尾不加标点的排法一致；译文按中文习惯仍以句号收尾。]

\yizhu{整节的路线可以概括为“先放大、再切掉”：借极丰富线丛把 $M$ 嵌进 $\PP^{N}$，使 $M$ 上的层都能用 $\cO(k)$ 的消解表示，于是 $F$ 的作用就转写成 $\PP^{N}\times X$ 上的对象 $\cE'$；再用 (36) 的卷积 $G$ 与同调维数估计把多余的部分 $C$ 分离出去，剩下 $\cE$。$K$ 与 $\cE'$ 的差别由特异三角 $S\to K\to\cE'\to S[1]$ 度量，而 $S$ 的上同调集中在 $[-n-a,-n]$，正好不与 $\Rf J_{*}\cE$ 所在的区间 $[-a,0]$ 相交。【来源：原书 p.549--555 全节结构及 p.555 末的论证。】}
